% 题干转写：完成
% 解析转写：完成
% 图形核对：完成
\section{第35届复赛第7题：弹簧衍射与DNA X射线衍射类比}
\subsection{题干}
用波长 $\lambda=633\,\mathrm{nm}$ 的激光水平照射竖直圆珠笔中的小弹簧，在距离弹簧 $L=4.2\,\mathrm m$ 的光屏上形成衍射图像，如图\ref{fig:cpho35q07}（a）。该图像与用波长 $\lambda'=0.15\,\mathrm{nm}$ 的平行X射线照射DNA样品、在距离 $L'=9.0\,\mathrm{cm}$ 的底片上获得的图\ref{fig:cpho35q07}（b）相似。

求：（1）利用图（a）的尺寸信息估算弹簧钢丝直径 $d_1$、弹簧半径 $R$ 和螺距 $p$；（2）利用图（b）的尺寸信息估算DNA螺旋半径 $R'$ 和螺距 $p'$。已知两段互成角度的细铁丝的衍射干涉图像，与两条成同样角度、相同宽度的狭缝图像一致。
\begin{figure}[htbp]\centering
\begin{tikzpicture}[>=Stealth,font=\small]
  % 图a：弹簧与交叉衍射条纹
  \begin{scope}
    \node[font=\bfseries] at (1.9,3.65) {（a）弹簧与激光衍射图样};
    \draw[thick] (0.5,0.4) -- (0.5,3.0);
    \draw[thick,domain=0:720,samples=180,smooth,variable=\t]
      plot ({1.05+0.30*sin(\t)}, {0.55+\t/300});
    \draw[thick] (1.55,0.4) -- (1.55,3.0);
    \node[below] at (1.03,0.35) {弹簧};
    \draw[fill=black] (3.35,0.45) rectangle (5.75,3.0);
    \foreach \yy in {0.72,0.98,1.24,1.50,1.76,2.02,2.28,2.54,2.80}{
      \draw[red!75!yellow,line width=1.2pt] (4.55,\yy) -- (3.55,{0.53+0.47*(\yy-0.55)}) -- (3.72,{0.53+0.47*(\yy-0.55)});
      \draw[red!75!yellow,line width=1.2pt] (4.55,\yy) -- (5.55,{0.53+0.47*(\yy-0.55)}) -- (5.38,{0.53+0.47*(\yy-0.55)});
      \draw[red!75!yellow,line width=1.2pt] (4.55,\yy) -- (3.55,{2.92-0.47*(\yy-0.55)}) -- (3.72,{2.92-0.47*(\yy-0.55)});
      \draw[red!75!yellow,line width=1.2pt] (4.55,\yy) -- (5.55,{2.92-0.47*(\yy-0.55)}) -- (5.38,{2.92-0.47*(\yy-0.55)});
    }
    \fill[red!80!yellow] (4.55,1.65) circle (0.12);
    \draw[<->] (5.95,0.45) -- (5.95,3.0) node[midway,right] {实际尺寸 $7.8\,\mathrm{cm}$};
    \node[below] at (4.55,0.25) {交叉衍射条纹};
  \end{scope}
  % 图b：DNA衍射图样
  \begin{scope}[xshift=7.4cm]
    \node[font=\bfseries] at (2.2,3.65) {（b）DNA 的 X 射线衍射图样};
    \fill[gray!12] (2.1,1.68) circle (1.48);
    \draw[thick] (2.1,1.68) circle (1.48);
    \fill[white] (2.1,1.68) circle (0.18);
    \draw (2.1,1.68) circle (0.18);
    \foreach \i/\yy in {1/0.42,2/0.68,3/0.94,5/1.20}{
      \fill[black] ({2.1-0.16*\i}, {1.68+\yy}) ellipse (0.14 and 0.055);
      \fill[black] ({2.1+0.16*\i}, {1.68+\yy}) ellipse (0.14 and 0.055);
      \fill[black] ({2.1-0.16*\i}, {1.68-\yy}) ellipse (0.14 and 0.055);
      \fill[black] ({2.1+0.16*\i}, {1.68-\yy}) ellipse (0.14 and 0.055);
    }
    \draw[densely dashed] (1.35,0.95) -- (2.85,2.41);
    \draw[densely dashed] (1.35,2.41) -- (2.85,0.95);
    \node[align=center,scale=0.8] at (2.1,0.02) {第4级亮纹缺失\\（示意）};
    \draw[<->] (3.9,0.2) -- (3.9,3.16) node[midway,right] {实际尺寸 $9.4\,\mathrm{cm}$};
  \end{scope}
\end{tikzpicture}

\caption{第35届复赛第7题图（a）、图（b）的结构性 TikZ 重绘。衍射斑纹的位置、夹角和尺寸基准用于表达题目测量信息，不以原始照片作为成品。}\label{fig:cpho35q07}
\end{figure}

\subsection{参考解答整理}
\paragraph{（1）弹簧的衍射测量。} 由图（a）测得两组条纹中点连线夹角
\[
\phi=30^\circ.
\]
远场小角近似下，第 $n$ 级暗纹满足
\[
d\sin\theta_n\simeq d\theta_n=d\frac{x_n}{L}=n\lambda,
\]
故相邻暗纹间距为
\[
X=x_{n+1}-x_n=\frac{\lambda L}{d}.
\]
图中长、短空间周期分别为 $X_1=7.9\,\mathrm{mm}$、$X_2=1.7\,\mathrm{mm}$，从而
\[
d_1=\frac{\lambda L}{X_1}=\boxed{0.34\,\mathrm{mm}},\qquad
d_2=\frac{\lambda L}{X_2}=\boxed{1.6\,\mathrm{mm}}.
\]
将螺旋线在平面展开，结合交叉细丝夹角，有
\[
p=\frac{d_2}{\cos(\phi/2)}
=\frac{\lambda L}{X_2\cos(\phi/2)}
=\boxed{1.6\,\mathrm{mm}},
\]
\[
R=\frac{d_2}{2\pi\sin(\phi/2)}
=\frac{\lambda L}{2\pi X_2\sin(\phi/2)}
=\boxed{0.96\,\mathrm{mm}}.
\]

\paragraph{（2）DNA螺旋参数。} 由图（b）量得
\[
\phi'=84^\circ,\qquad X'=4.7\,\mathrm{mm}.
\]
先由衍射周期得到相邻分子链间距
\[
d'=\frac{\lambda'L'}{X'}=\boxed{2.9\,\mathrm{nm}}.
\]
套用同一螺旋展开几何关系，
\[
p'=\frac{d'}{\cos(\phi'/2)}
=\frac{\lambda'L'}{X'\cos(\phi'/2)}
=\boxed{3.9\,\mathrm{nm}},
\]
\[
R'=\frac{d'}{2\pi\sin(\phi'/2)}
=\frac{\lambda'L'}{2\pi X'\sin(\phi'/2)}
=\boxed{0.68\,\mathrm{nm}}.
\]

\subsection{转写依据}
题干、答案公式和图件分别依据 \texttt{CPHO35-Q07.md}、\texttt{CPHO35-Q07\_formula\_check\_v1.csv} 及图件核对入口；当前PDF未明示分值，未写入图册。
